% Nouveau EGA V (cote 25): archive cote 25, site page 144 (PDF page 145).
% typescript. Draft compared: none.
% AI-transcribed by Claude Opus 5.5 (claude-opus-5-5, adaptive thinking, effort high), Message Batches API, batch ega25b (msgbatch_01A7n5YuXubUPNK1zsepi8RG), 2026-11-13; not yet checked by eye.
\scanpage{25/144}
\mspage{\hand{5}}

La colonne provient de la suite exacte de transtitivité pour les
morphismes lisses $X\hand{\to} Y$ et $Y\hand{\to} S$ , en  tensorisant par $\underline{O}_{V1}$ (cela reste
exact, car tous les modules de la suite sont loc libres). La ligne
horizontale est un cas particulier d'une suite exacte obtenue chaque
fois que sur X  sur S, on a une section \hand{$\mu$} d'un module loc libre F
et qu'on prend le schéma W des zéros; on trouve une suite exacten
\[ \overset{\hand{\vee}}{F}\otimes \underline{O}_X \hand{\longrightarrow} \underline{\hand{\Omega}}^1_{X/S}\otimes \underline{O}_W \hand{\longrightarrow} \underline{\hand{\Omega}}^1_{W/S} \hand{\longrightarrow} 0 \]
et (si X/S est lisse) le premier homomorphisme est injectif exactement en les points où
W est lisse sur X avec la ``bonne '' \unsure{dimension} relative (i.e. partout dans

le cas actuel).  Cette suite exacte est 
  une
conséquence immédiate de la suite exacte
\[ \underline{J}/\underline{J}^2\hand{\longrightarrow} \underline{\hand{\Omega}}^1_{X/S}\otimes \underline{O}_W \hand{\longrightarrow} \underline{\hand{\Omega}}^1_{X/S} \hand{\longrightarrow} 0 \]
\placedmargin{16 \textbar} qui figure au par.16 (on y pourrait mettre en corollaire la version
signalée ici). La caractérisation de l'ensemble des points où on
peut mettre un zéro à gauche est contenue dans le critère Jacobien.

Notons que l'on a un isomorphisme canonique
\[ \underline{P}^1_{X/S} = \underline{\hand{\Omega}}^1_{X/S} + \underline{O}_X \ ,\ \text{d'où}\ \overset{\hand{\scriptscriptstyle\vee}}{\underline{P}}^1_{X/S} = \hand{\mathcal{T}}_{X/S} + \underline{O}_X \ , \]
d'autre part on \unsure{vérifie|vérifié} que l'homomorphisme composé \hand{$\mu$} du diagramme \hand{\unsure{(D)|(1)}} est nul
sur le facteur $\underline{O}_{V1}$ , et sur le facteur $\hand{\mathcal{T}}_{X/S}\otimes\underline{O}_{V1}$ il se réduit à
l'homomorphisme $M(\hand{\phi})'$ déduit de la section $M(\hand{\phi})$ de $\underline{\mathrm{Sym}}^2(\hand{\underline{\hand{\Omega}}}^1_{X/S})\otimes \underline{O}_{V1}$ déjà
envisagée. Donc en un point x de X, $M(\hand{\phi})$ est non dégénéré i.e. $M(\hand{\phi})'$
est surjectif sss \hand{$\mu$} est surjectif en x, et on voit sur le diagramme \hand{(D)},
que cela équivaut aussi à dire que \hand{$\nu$} est surjectif en x (car l'un et l'autre
signifient que l'homomorphisme canonique de la somme des deux sous
Modules envisagés de $\underline{\hand{\Omega}}^1_{X/S}\otimes \underline{O}_{V1}$ dans ce derniers est surjectif en
x). On trouve ainsi :

\mspage{\hand{144}}
